\section{Related Work}
\label{sec:related}

\subsection{Honesty, accuracy, and operational belief}

\para{Behavioral definitions.} \citet{ren2025mask} define model honesty by comparing pressured statements with beliefs elicited under neutral conditions. This separates whether a report matches an elicited belief from whether that belief matches ground truth. \citet{kretschmar2025liars} extend the evaluation surface across seven settings and explicitly show why factual correctness is not a lie label: a model can honestly state a wrong belief or report a true statement that contradicts its belief. Our study adopts this distinction but focuses on a narrower, matched factual-QA panel in which lie and error mechanisms can be compared for one model and one question distribution.

\para{Limits of belief verification.} A repeated answer is evidence of behavioral stability, not direct access to a mental state. \citet{cooney2026did} make this concern concrete by evaluating belief-verified model organisms; activation and log-probability detectors that work on prompted lies degrade on these harder cases. The Anthropic Fellows study similarly reports label revisions after stronger adjudication and weak cross-category transfer \citep{hopkins2026finetuned}. We therefore use ``belief'' only as an operational shorthand, preserve ambiguous pressured errors rather than forcing binary labels, and avoid claims about awareness or intent.

\subsection{White-box lie and truth probes}

Early activation classifiers separated true from false statements using hidden states \citep{azaria2023internal}. \citet{burger2024truth} later proposed a polarity-robust truth subspace and found broadly shared truth directions. These methods establish that internal activations can predict factual properties, but they do not by themselves establish a representation of belief-contradicting intent. Likewise, the unrelated-question detector of \citet{pacchiardi2023catch} suggests that behavioral traces can transfer across some black-box lying settings, yet later benchmarks show that transfer is strongly setting-dependent.

The closest methodological baseline is the residual-stream logistic probe of \citet{goldowskydill2025strategic}. It attains 0.96--0.999 AUROC in its evaluation settings, while its documented false triggers include honest text about deception and requests for incorrect content. Broader studies sharpen this caveat. \citet{natarajan2026one} find that matching probes to known deception types substantially improves performance and that the training instruction pair explains much of the variation. \citet{moustafa2026beyond} show that optimal depth and probe benefit depend on the dataset and lie typology. \citet{luikham2026asymmetries} find only partial geometric overlap and asymmetric transfer between spontaneous and instructed deception. Finally, \citet{thormann2026limits} shows that truth probes can miss misleading behavior that contains no false statement. Building on this work, we ask a more controlled question: does a standard probe distinguish two false-report mechanisms when prompt context is matched?

\subsection{Hallucination and uncertainty detection}

\citet{kuhn2023semantic} estimate semantic uncertainty by clustering sampled generations according to meaning and computing entropy over the induced classes. This is well suited to unstable answers but need not identify a confident wrong association or a confident lie. \citet{cheang2025recall} formalize a related distinction between unassociated and associated hallucinations and find that hidden-state detectors work much better on weakly grounded errors than on strongly associated wrong answers. These results motivate our three-way neutral stratification: stable correct, stable wrong, and unstable.

Our short-answer entropy score is an intentionally simple approximation to semantic entropy. Because the same five samples determine both entropy and stability labels, it is a manipulation check rather than an independent test of detector generalization. We complement it with a correctness probe trained only on a disjoint honest-request set. This combination lets us test whether a putative lie detector behaves more like a prompt classifier, a correctness detector, or a mechanism-specific monitor.
